% Pandoc LaTeX header include for jobsheet PDFs (lualatex).

% Emoji fallback so callout markers (checkpoint, warning) render instead of
% showing as tofu boxes: DejaVu Sans (the body font) doesn't cover emoji.
% Map each emoji codepoint used in the jobsheets to a glyph pulled from a
% font that does cover it, via newunicodechar (more reliable under lualatex
% than a blanket fontspec fallback feature).
%
% Font choice varies by machine, so fall back down a chain instead of
% hardcoding one: Segoe UI Emoji (present on this dev machine) -> Symbola
% (monochrome, installed via `brew install --cask font-symbola` on the
% self-hosted macOS CI runner -- avoid Apple Color Emoji, its sbix color
% glyphs render poorly under lualatex) -> DejaVu Sans as a last resort
% (covers warning-sign ⚠ but not the checkmark ✅, which would tofu there).
\usepackage{fontspec}
\IfFontExistsTF{Segoe UI Emoji}
  {\newfontfamily{\emojifont}{Segoe UI Emoji}}
  {\IfFontExistsTF{Symbola}
    {\newfontfamily{\emojifont}{Symbola}}
    {\newfontfamily{\emojifont}{DejaVu Sans}}}
\usepackage{newunicodechar}
\newunicodechar{✅}{{\emojifont ✅}}
\newunicodechar{⚠}{{\emojifont ⚠}}
% U+FE0F (variation selector-16, often trailing "⚠️") has no visible glyph
% of its own — drop it instead of leaving it unmapped.
\newunicodechar{️}{}

% Style blockquotes (used for "Checkpoint" / "Jika gagal" callouts) as a
% light-gray left-rule box instead of plain indented italics. \FrameCommand
% is redefined only inside a local group scoped to \quote, so it doesn't
% leak into pandoc's own Shaded/snugshade environment used for code blocks
% (both are built on the framed package and share that command globally).
\usepackage{framed}
\usepackage{xcolor}
\definecolor{calloutrule}{HTML}{94A3B8}
\definecolor{calloutbg}{HTML}{F8FAFC}
\renewenvironment{quote}
  {\begingroup
   \def\FrameCommand{{\color{calloutrule}\vrule width 3pt}\colorbox{calloutbg}}
   \setlength{\FrameSep}{8pt}
   \begin{leftbar}}
  {\end{leftbar}\endgroup}

% Keep code blocks from overflowing the page width. fvset (rather than
% redefining the Highlighting environment outright) preserves pandoc's own
% commandchars setup, which is what lets syntax-highlighting token commands
% inside code blocks be interpreted instead of printed literally.
\usepackage{fvextra}
\fvset{breaklines=true, breakanywhere=true, fontsize=\small}
